% ==========================================================================
%  Chapter: Maxwell's equations in vacuum
%  Conventions: Gaussian units; metric signature (+,-,-,-); x^0 = ct.
%  Packages needed in the master preamble: amsmath, amssymb, amsthm,
%  booktabs, tikz.  Theorem-like environments used: theorem, proposition,
%  definition, remark, example, problem (defined below if missing).
% ==========================================================================

\makeatletter
\@ifundefined{theorem}{\newtheorem{theorem}{Theorem}[chapter]}{}
\@ifundefined{proposition}{\newtheorem{proposition}{Proposition}[chapter]}{}
\theoremstyle{definition}
\@ifundefined{definition}{\newtheorem{definition}{Definition}[chapter]}{}
\@ifundefined{example}{\newtheorem{example}{Example}[chapter]}{}
\@ifundefined{problem}{\newtheorem{problem}{Problem}[chapter]}{}
\theoremstyle{remark}
\@ifundefined{remark}{\newtheorem{remark}{Remark}[chapter]}{}
\theoremstyle{plain}
\makeatother

\providecommand{\vE}{\mathbf{E}}
\providecommand{\vB}{\mathbf{B}}
\providecommand{\vA}{\mathbf{A}}
\providecommand{\vj}{\mathbf{j}}
\providecommand{\vk}{\mathbf{k}}
\providecommand{\vx}{\mathbf{x}}
\providecommand{\vn}{\mathbf{n}}
\providecommand{\vS}{\mathbf{S}}
\providecommand{\vg}{\mathbf{g}}
\providecommand{\vv}{\mathbf{v}}
\providecommand{\vK}{\mathbf{K}}
\providecommand{\vF}{\mathbf{F}}
\providecommand{\vf}{\mathbf{f}}
\providecommand{\vy}{\mathbf{y}}
\providecommand{\vbeta}{\boldsymbol{\beta}}
\providecommand{\grad}{\boldsymbol{\nabla}}
\providecommand{\curl}{\boldsymbol{\nabla}\!\times}
\providecommand{\divg}{\boldsymbol{\nabla}\!\cdot}
\providecommand{\dd}{\mathrm{d}}
\providecommand{\ii}{\mathrm{i}}
\providecommand{\pa}{\partial}
\providecommand{\Ftil}{\tilde F}

\chapter{Maxwell's Equations in Vacuum}
\label{ch:maxwell-vacuum}

The Maxwell equations are the oldest field theory that survived the
arrival of relativity without a single change of content, and for that
reason they are the natural laboratory for every general notion of field
theory: constraints, gauge freedom, counting of degrees of freedom,
well-posedness of initial and boundary problems, causal structure,
conformal and duality symmetry.  In this chapter we treat them in full,
in the simplest setting---empty space, with the charges and currents
regarded as given sources---and we try to keep three questions apart
that are often blurred in textbooks:
\begin{enumerate}
\item \emph{What do the equations say?}\ \ (Sections~\ref{sec:mx-eqs}--\ref{sec:mx-geom}):
  the equations in three-dimensional and covariant form, their
  variational origin, and their geometry.
\item \emph{How much freedom does a solution have?}\ \ (Section~\ref{sec:mx-dof}):
  the number of degrees of freedom, counted in several independent ways.
\item \emph{What data single out a solution?}\ \ (Sections~\ref{sec:mx-free}--\ref{sec:mx-num}):
  initial-value, source, boundary, exterior and other problems, and the
  question of when each is well posed.
\end{enumerate}
Along the way we flag, in remarks, the places where an object of the
\emph{ideal world} of the model (unbounded space, point source, perfect
conductor, field without sources) is silently identified with
an object of the \emph{real world} of experiments.

% ==========================================================================
\section{Scope and conventions}\label{sec:mx-conv}
% ==========================================================================

\subsection{What ``in vacuum'' means}

``In vacuum'' means that the only fields are $\vE$ and $\vB$; there are no
auxiliary fields $\mathbf D,\mathbf H$, no polarization or magnetization,
no constitutive relations.  Charges and currents may be present, but they
are the \emph{microscopic} ones: a charge density $\rho(t,\vx)$ and a
current density $\vj(t,\vx)$, given functions (or distributions) that
satisfy the continuity equation.  The equations are then the
Maxwell--Lorentz equations.  The special case $\rho=0,\ \vj=0$ is the
\emph{free} (source-free) field; it is of special importance, because the
general solution is a particular solution with sources plus a free
field, and the free fields carry all the dynamical information (the
``degrees of freedom'') of the theory.

\subsection{Conventions}

\begin{itemize}
\item \textbf{Units.}  Gaussian: $\vE$ and $\vB$ have the same dimension,
  Coulomb's law reads $F=q_1q_2/r^2$, the Lorentz force on a point charge
  is $\vF=q\,(\vE+\vv\times\vB/c)$, and $c$ appears explicitly.  The
  dictionary to SI is
  \begin{center}
  \begin{tabular}{@{}ccccccc@{}}
  \toprule
  $\vE_{\mathrm G}$ & $\vB_{\mathrm G}$ & $\rho_{\mathrm G}$ & $\vj_{\mathrm G}$ & $\varphi_{\mathrm G}$ & $\vA_{\mathrm G}$ & $c^2$\\
  \midrule
  $\sqrt{4\pi\varepsilon_0}\,\vE_{\mathrm{SI}}$ &
  $\sqrt{4\pi/\mu_0}\,\vB_{\mathrm{SI}}$ &
  $\rho_{\mathrm{SI}}/\sqrt{4\pi\varepsilon_0}$ &
  $\vj_{\mathrm{SI}}/\sqrt{4\pi\varepsilon_0}$ &
  $\sqrt{4\pi\varepsilon_0}\,\varphi_{\mathrm{SI}}$ &
  $\sqrt{4\pi/\mu_0}\,\vA_{\mathrm{SI}}$ &
  $1/\mu_0\varepsilon_0$\\
  \bottomrule
  \end{tabular}
  \end{center}
\item \textbf{Metric.}  $\eta_{\mu\nu}=\eta^{\mu\nu}=\operatorname{diag}(+1,-1,-1,-1)$;
  $x^\mu=(ct,\vx)$, $x_\mu=(ct,-\vx)$.  Greek indices run over $0,1,2,3$,
  Latin ones over $1,2,3$; repeated indices are summed.  Spatial
  vectors have components $V^i$ and are \emph{not} distinguished from
  $V_i$ when written in three-dimensional notation ($V_i\equiv V^i$).
\item \textbf{Derivatives.}  $\pa_\mu=\pa/\pa x^\mu=(c^{-1}\pa_t,\grad)$,
  $\pa^\mu=(c^{-1}\pa_t,-\grad)$, and
  $\Box=\pa_\mu\pa^\mu=c^{-2}\pa_t^2-\nabla^2$.
\item \textbf{Levi-Civita symbol.}  $\varepsilon^{0123}=+1$, hence
  $\varepsilon_{0123}=-1$ (because $\det\eta=-1$).
\item \textbf{Sources and potentials.}  $j^\mu=(c\rho,\vj)$ and
  $A^\mu=(\varphi,\vA)$.
\end{itemize}

% ==========================================================================
\section{The equations in three-dimensional form}\label{sec:mx-eqs}
% ==========================================================================

\subsection{The four equations and the force law}

\begin{subequations}\label{eq:maxwell-3d}
\begin{align}
\divg\vE &= 4\pi\rho, \label{eq:gauss-E}\\
\divg\vB &= 0, \label{eq:gauss-B}\\
\curl\vE + \frac1c\,\frac{\pa\vB}{\pa t} &= 0, \label{eq:faraday}\\
\curl\vB - \frac1c\,\frac{\pa\vE}{\pa t} &= \frac{4\pi}{c}\,\vj. \label{eq:ampere}
\end{align}
\end{subequations}
Equations \eqref{eq:gauss-B} and \eqref{eq:faraday} are \emph{homogeneous}
(they do not contain the sources), while \eqref{eq:gauss-E} and
\eqref{eq:ampere} are \emph{inhomogeneous}.  The field acts back on the
charges through the Lorentz force; its density is
\begin{equation}\label{eq:force-density}
\vf=\rho\,\vE+\frac1c\,\vj\times\vB ,
\end{equation}
and for a point charge $q$ with velocity $\vv$ it gives
$\vF=q(\vE+\vv\times\vB/c)$.

\begin{remark}[Real world vs.\ ideal world: what $\vE$ and $\vB$ are]
The fields are not ``seen''; operationally they are \emph{defined} by
the force law: $\vE$ and $\vB$ at a point are the vectors that, inserted
into $q(\vE+\vv\times\vB/c)$, reproduce the measured force on a small test
charge. Causality and experimental support belong to this operational
level.  The partial differential equations \eqref{eq:maxwell-3d} are a
statement in the ideal world of the model: they speak about fields at
every point of an idealized continuum and for all times, including
regions and moments where no test charge was ever placed.  Maxwell's
equations without the force law are a mathematical theory without a
physical referent; the force law without the equations is a rule
without dynamics.  Both are needed, and in what follows we shall always
regard $\rho,\vj$ as given when we solve for the fields, even though in
reality the force law couples them back.
\end{remark}

\subsection{Structure: constraints and evolution equations}

The four equations are not on the same footing.  Equations
\eqref{eq:gauss-E} and \eqref{eq:gauss-B} contain no time derivatives:
they are \emph{constraints} on the field at a given instant.
Equations \eqref{eq:faraday} and \eqref{eq:ampere} are first-order
\emph{evolution equations}: they give $\pa_t\vB$ and $\pa_t\vE$ in terms
of the instantaneous fields,
\begin{equation}\label{eq:evolution}
\frac{\pa\vE}{\pa t}=c\,\curl\vB-4\pi\vj,\qquad
\frac{\pa\vB}{\pa t}=-c\,\curl\vE .
\end{equation}
The central structural fact is that the evolution equations
\emph{propagate} the constraints.  Taking the divergence of
\eqref{eq:faraday} and using $\divg\curl\equiv0$,
\begin{equation}\label{eq:propagate-B}
\frac{\pa}{\pa t}\,(\divg\vB)=0 ,
\end{equation}
and taking the divergence of \eqref{eq:ampere},
\begin{equation}\label{eq:propagate-E}
\frac{\pa}{\pa t}\,(\divg\vE-4\pi\rho)
=-4\pi\Big(\frac{\pa\rho}{\pa t}+\divg\vj\Big).
\end{equation}
Therefore: (i) if the constraints hold at $t=0$ they hold for all $t$, and
(ii) the \emph{system is consistent only if the sources satisfy the
continuity equation}
\begin{equation}\label{eq:continuity}
\frac{\pa\rho}{\pa t}+\divg\vj=0\qquad\Longleftrightarrow\qquad
\pa_\mu j^\mu=0 .
\end{equation}
Historically this is the reason for Maxwell's displacement current
$-c^{-1}\pa_t\vE$ in \eqref{eq:ampere}: without it the divergence of
Ampère's law would force $\divg\vj=0$ for all time-dependent charge
distributions.  Charge conservation is thus not an extra postulate but a
\emph{compatibility condition} of the field equations.

\subsection{Integral form and jump conditions}

For a fixed volume $V$ with boundary $\pa V$ and a fixed surface $\Sigma$
with boundary $\pa\Sigma$, \eqref{eq:maxwell-3d} is equivalent (for
sufficiently regular fields) to
\begin{align}
\oint_{\pa V}\vE\cdot\dd\mathbf a&=4\pi Q_V,&
\oint_{\pa V}\vB\cdot\dd\mathbf a&=0,\\
\oint_{\pa\Sigma}\vE\cdot\dd\boldsymbol\ell&=-\frac1c\frac{\dd}{\dd t}\int_\Sigma\vB\cdot\dd\mathbf a,&
\oint_{\pa\Sigma}\vB\cdot\dd\boldsymbol\ell&=\frac1c\frac{\dd}{\dd t}\int_\Sigma\vE\cdot\dd\mathbf a+\frac{4\pi}{c}\int_\Sigma\vj\cdot\dd\mathbf a .
\end{align}
The integral form is the more fundamental one when the fields are not
differentiable.  Applied to thin pillboxes and loops across a surface
$\mathcal S$ with unit normal $\vn$ that carries a surface charge density
$\sigma$ and a surface current density $\vK$, it yields the jump
conditions
\begin{equation}\label{eq:jump}
\vn\cdot[\vE]=4\pi\sigma,\quad
\vn\times[\vE]=0,\quad
\vn\cdot[\vB]=0,\quad
\vn\times[\vB]=\frac{4\pi}{c}\vK ,
\end{equation}
where $[\,\cdot\,]$ is the value on the side into which $\vn$ points minus
the value on the other side.  These will be the boundary conditions of
Section~\ref{sec:mx-bvp}.

\subsection{Energy, momentum and angular momentum}

Define the energy density, the Poynting vector, the momentum density, and
the Maxwell stress tensor
\begin{equation}\label{eq:densities}
u=\frac{\vE^2+\vB^2}{8\pi},\qquad
\vS=\frac{c}{4\pi}\,\vE\times\vB,\qquad
\vg=\frac{\vS}{c^2},\qquad
\sigma_{ij}=\frac1{4\pi}\Big[E_iE_j+B_iB_j-\tfrac12\delta_{ij}(\vE^2+\vB^2)\Big].
\end{equation}
From \eqref{eq:maxwell-3d} one derives the \emph{Poynting theorem} (using the identity
$\divg(\vE\times\vB)=\vB\cdot\curl\vE-\vE\cdot\curl\vB$) and the
\emph{momentum balance}
\begin{equation}\label{eq:poynting}
\frac{\pa u}{\pa t}+\divg\vS=-\vj\cdot\vE,\qquad
\frac{\pa g_i}{\pa t}-\pa_j\sigma_{ij}=-f_i ,
\end{equation}
with $\vf$ from \eqref{eq:force-density}.  The right-hand sides are the
work done by the field on the charges and the force exerted on them: in
the source-free case $u,\vS$ and $\vg$, $\sigma_{ij}$ obey exact local
conservation laws.  The angular-momentum density is $\vx\times\vg$, and
it obeys a conservation law (without sources) because $\sigma_{ij}$ is
symmetric.  The total energy, momentum and angular momentum
\begin{equation}
W=\int u\,\dd^3x,\qquad \mathbf P=\int\vg\,\dd^3x,\qquad
\mathbf J=\int \vx\times\vg\,\dd^3x
\end{equation}
are conserved for the free field (with appropriate fall-off at infinity).
The positivity $u\ge0$ and the inequality
\begin{equation}\label{eq:S-bound}
|\vS|\le c\,u ,
\end{equation}
which follows from $|\vE\times\vB|\le|\vE||\vB|\le\tfrac12(\vE^2+\vB^2)$,
are the two elementary facts from which well-posedness and finite
propagation speed will be derived in Section~\ref{sec:mx-cauchy}.

% ==========================================================================
\section{Covariant formulation and the action}\label{sec:mx-cov}
% ==========================================================================

\subsection{The field-strength tensor}

The six components of $\vE$ and $\vB$ are assembled into the
antisymmetric tensor $F_{\mu\nu}$.  With $A_\mu=(\varphi,-\vA)$,
\begin{equation}\label{eq:potentials}
F_{\mu\nu}=\pa_\mu A_\nu-\pa_\nu A_\mu ,\qquad
\vE=-\grad\varphi-\frac1c\frac{\pa\vA}{\pa t},\qquad
\vB=\curl\vA ,
\end{equation}
so that $F_{0i}=E_i$ and $F_{ij}=-\varepsilon_{ijk}B_k$.  In matrix form
(row index $\mu$, column index $\nu$):
\begin{equation}\label{eq:Fmatrix}
F^{\mu\nu}=\begin{pmatrix}
0&-E_x&-E_y&-E_z\\
E_x&0&-B_z&B_y\\
E_y&B_z&0&-B_x\\
E_z&-B_y&B_x&0
\end{pmatrix},\qquad
\Ftil^{\mu\nu}\equiv\tfrac12\varepsilon^{\mu\nu\rho\sigma}F_{\rho\sigma}=\begin{pmatrix}
0&-B_x&-B_y&-B_z\\
B_x&0&E_z&-E_y\\
B_y&-E_z&0&E_x\\
B_z&E_y&-E_x&0
\end{pmatrix}.
\end{equation}
The dual tensor $\Ftil^{\mu\nu}$ is obtained from $F^{\mu\nu}$ by the
\emph{duality substitution} $\vE\to\vB,\ \vB\to-\vE$.

Relative to an observer with four-velocity $u^\mu$, $u_\mu u^\mu=1$
(dimensionless normalization), the electric and magnetic fields are the
four-vectors
\begin{equation}\label{eq:EB-observer}
E^\mu=F^{\mu\nu}u_\nu,\qquad B^\mu=\Ftil^{\mu\nu}u_\nu,\qquad
F^{\mu\nu}=E^\mu u^\nu-E^\nu u^\mu+\varepsilon^{\mu\nu\rho\sigma}u_\rho B_\sigma .
\end{equation}
For $u^\mu=(1,\mathbf 0)$ these are $E^\mu=(0,\vE)$, $B^\mu=(0,\vB)$.  The
splitting into ``electric'' and ``magnetic'' is therefore
observer-dependent: only $F_{\mu\nu}$ is a frame-independent object.
Under a boost to a frame moving with velocity $c\vbeta$ relative to the
original one (with $\gamma=(1-\beta^2)^{-1/2}$),
\begin{equation}\label{eq:boost}
\vE'_\parallel=\vE_\parallel,\quad \vB'_\parallel=\vB_\parallel,\qquad
\vE'_\perp=\gamma(\vE+\vbeta\times\vB)_\perp,\quad
\vB'_\perp=\gamma(\vB-\vbeta\times\vE)_\perp ,
\end{equation}
where $\parallel,\perp$ refer to $\vbeta$.  A purely electric field in
one frame is therefore magnetic in another; and the Lorentz force law
takes the covariant form $\dd p^\mu/\dd\tau=(e/c)\,F^{\mu\nu}u_\nu$ (here
$u^\nu=\dd x^\nu/\dd\tau$ with $u^2=c^2$).

\subsection{Equations of motion}

With $j^\mu=(c\rho,\vj)$ the four equations \eqref{eq:maxwell-3d} collapse
into two:
\begin{equation}\label{eq:maxwell-cov}
\pa_\mu F^{\mu\nu}=\frac{4\pi}{c}\,j^\nu,\qquad
\pa_\mu\Ftil^{\mu\nu}=0
\quad\Longleftrightarrow\quad
\pa_\lambda F_{\mu\nu}+\pa_\mu F_{\nu\lambda}+\pa_\nu F_{\lambda\mu}=0 .
\end{equation}
(For $\nu=0$ the first gives \eqref{eq:gauss-E}, for $\nu=i$ it gives
\eqref{eq:ampere}; the second gives \eqref{eq:gauss-B} and
\eqref{eq:faraday}.)  The second equation is the \emph{Bianchi identity};
it is the integrability condition for the existence of $A_\mu$ and
holds identically once $F_{\mu\nu}=\pa_\mu A_\nu-\pa_\nu A_\mu$.  The
continuity equation follows from the first by antisymmetry:
$\pa_\nu\pa_\mu F^{\mu\nu}\equiv0\Rightarrow\pa_\nu j^\nu=0$.

\subsection{Invariants}

There are two independent quadratic Lorentz invariants,
\begin{equation}\label{eq:invariants}
F_{\mu\nu}F^{\mu\nu}=2(\vB^2-\vE^2),\qquad
F_{\mu\nu}\Ftil^{\mu\nu}=-4\,\vE\cdot\vB .
\end{equation}
The first is a scalar, the second a pseudoscalar.

\subsection{Gauge freedom and the Lorenz gauge}

The potential is not determined by $F_{\mu\nu}$: for any function
$\chi(x)$,
\begin{equation}\label{eq:gauge}
A^\mu\to A^\mu+\pa^\mu\chi\qquad\Longleftrightarrow\qquad
\varphi\to\varphi+\frac1c\,\pa_t\chi,\quad \vA\to\vA-\grad\chi,
\end{equation}
leaves $\vE$ and $\vB$ unchanged.  The first equation of
\eqref{eq:maxwell-cov} becomes $\Box A^\nu-\pa^\nu(\pa_\mu A^\mu)=\tfrac{4\pi}{c}j^\nu$.
In the \emph{Lorenz gauge} $\pa_\mu A^\mu=c^{-1}\pa_t\varphi+\divg\vA=0$,
\begin{equation}\label{eq:lorenz-wave}
\Box A^\mu=\frac{4\pi}{c}\,j^\mu,\qquad \pa_\mu A^\mu=0 ,
\end{equation}
four decoupled wave equations.  A residual gauge freedom $\Box\chi=0$
remains.

\subsection{Stress--energy tensor}

The symmetric, gauge-invariant, traceless tensor
\begin{equation}\label{eq:T}
T^{\mu\nu}=\frac1{4\pi}\Big(-F^{\mu\alpha}F^{\nu}{}_{\alpha}+\tfrac14\,\eta^{\mu\nu}F_{\alpha\beta}F^{\alpha\beta}\Big),\qquad
T^\mu{}_\mu=0,
\end{equation}
has components $T^{00}=u$, $T^{0i}=S^i/c$, $T^{ij}=-\sigma_{ij}$, and
satisfies
\begin{equation}\label{eq:T-conservation}
\pa_\mu T^{\mu\nu}=-\frac1c\,F^{\nu\lambda}j_\lambda\equiv-f^\nu,\qquad
f^\mu=\Big(\frac{\vj\cdot\vE}{c},\ \rho\vE+\frac1c\vj\times\vB\Big),
\end{equation}
which contains both equations of \eqref{eq:poynting}.  Tracelessness
reflects the absence of any length scale in the theory and is the origin
of its conformal invariance (Section~\ref{sec:mx-geom}).

\subsection{Action principle}

Let $\dd^4x=c\,\dd t\,\dd^3x$.  The field equations are the
Euler--Lagrange equations of
\begin{equation}\label{eq:action}
S=-\frac1{16\pi c}\int F_{\mu\nu}F^{\mu\nu}\,\dd^4x-\frac1{c^2}\int j^\mu A_\mu\,\dd^4x
=\int\dd t\,\dd^3x\ \mathcal L,\qquad
\mathcal L=\frac{\vE^2-\vB^2}{8\pi}-\rho\varphi+\frac1c\,\vj\cdot\vA .
\end{equation}
Indeed, $\pa\mathcal L/\pa(\pa_\mu A_\nu)=-F^{\mu\nu}/4\pi$ and
$\pa\mathcal L/\pa A_\nu=-j^\nu/c$, so
$\pa_\mu\big[\pa\mathcal L/\pa(\pa_\mu A_\nu)\big]=\pa\mathcal L/\pa A_\nu$
reads $\pa_\mu F^{\mu\nu}=(4\pi/c)j^\nu$; the homogeneous equations are
identities of the ansatz $F=\dd A$.  Under \eqref{eq:gauge} the action
changes by $\delta S=-c^{-2}\int j^\mu\pa_\mu\chi\,\dd^4x
=c^{-2}\int\chi\,\pa_\mu j^\mu\,\dd^4x+\text{boundary term}$, which vanishes
when the current is conserved.  Gauge invariance of the action and
conservation of the source are the same statement; this is the first
(and simplest) instance of Noether's second theorem:
$\pa_\nu\big(\delta S/\delta A_\nu\big)\equiv0$ is an \emph{identity} among
the Euler--Lagrange expressions, so the four field equations are not
independent.

% ==========================================================================
\section{Geometry of the field}\label{sec:mx-geom}
% ==========================================================================

\subsection{Minkowski space and the light cone}

The arena is Minkowski space $(\mathbb R^{1,3},\eta)$.  The set
$x^2=c^2t^2-|\vx|^2=0$ is the light cone; points with $x^2>0$ are
timelike, with $x^2<0$ spacelike, separated from the origin.  The
\emph{causal structure}---who can influence whom---is defined by the
cones, and the Poincaré group is the group of transformations that
preserve $\eta$.  We shall see that the Maxwell equations are
\emph{hyperbolic} with characteristics exactly the light cones, so that
disturbances propagate on and inside the cone, and that this single
geometric fact is behind uniqueness, finite propagation speed and
Huygens' principle.

\subsection{Differential-form language}

Since $F_{\mu\nu}$ is antisymmetric, it is naturally a 2-form,
\begin{equation}
\begin{aligned}
F&=\tfrac12F_{\mu\nu}\,\dd x^\mu\wedge\dd x^\nu
=E_i\,\dd x^0\wedge\dd x^i-\tfrac12\varepsilon_{ijk}B_k\,\dd x^i\wedge\dd x^j,\\
A&=A_\mu\,\dd x^\mu=\varphi\,\dd x^0-\vA\cdot\dd\vx,\qquad F=\dd A .
\end{aligned}
\end{equation}
Let $\star$ be the Hodge operator of $\eta$ with the orientation
$\varepsilon_{0123}=-1$ and define $(\star F)_{\mu\nu}=\tfrac12\varepsilon_{\mu\nu\rho\sigma}F^{\rho\sigma}=\Ftil_{\mu\nu}$ and
$(\star J)_{\mu\nu\rho}=\varepsilon_{\mu\nu\rho\sigma}j^\sigma$.  The Maxwell
equations are then simply
\begin{equation}\label{eq:forms}
\dd F=0,\qquad \dd{\star F}=\frac{4\pi}{c}\,{\star J},\qquad\text{hence}\qquad \dd{\star J}=0 .
\end{equation}
(One checks directly: $(\dd{\star F})_{123}=\divg\vE$ and
$(\star J)_{123}=c\rho$, so that \eqref{eq:gauss-E} is the $123$ component;
$(\dd{\star F})_{012}=c^{-1}\pa_tE_z-(\curl\vB)_z$ and
$(\star J)_{012}=-j_z$ give the $z$-component of \eqref{eq:ampere}.)  The
metric enters only through $\star$, so the only metric information
used by the dynamics is the \emph{conformal class} of $\eta$ (see
below).  Equations \eqref{eq:forms} make the following geometric facts
transparent:
\begin{itemize}
\item $\dd F=0$ is a statement of \emph{topology}: it makes no
  reference to the metric or to sources, and by Stokes' theorem it says
  that $\int_{\pa\mathcal C}F=0$ for the boundary of every
  three-dimensional region $\mathcal C$ of spacetime; this packages
  Faraday's law and $\oint\vB\cdot\dd\mathbf a=0$.
\item $\dd{\star F}=(4\pi/c)\star J$ is the only equation that contains
  dynamics (the metric) and sources.
\item $\dd{\star J}=0$ is charge conservation, obtained from $\dd^2=0$.
\item By the Poincaré lemma, $\dd F=0$ implies $F=\dd A$ \emph{locally}, and
  $A$ is determined up to $A\to A+\dd\chi$: this is the gauge freedom.
  Globally, $F=\dd A$ requires that the periods of $F$ over all closed
  2-cycles vanish, so the space of closed-but-not-exact fields is the
  second de Rham cohomology $H^2$ of the spacetime manifold.
\end{itemize}

\begin{remark}[The potential as a connection]
Geometrically $A$ is a connection on a principal $U(1)$ bundle and $F$
its curvature; a charged particle with charge $e$ transported around a
closed loop $\gamma$ acquires the phase factor $\exp\!\big(\tfrac{\ii e}{\hbar c}\oint_\gamma A\big)$.  On a space with a non-contractible loop (a long solenoid
region excluded from space) $F=0$ outside, yet $\oint A\neq0$: the
gauge-invariant content of $A$ is not only $F$ but the holonomies.  In
classical electrodynamics only $F$ is observable, and the ``gauge
freedom'' is the absence of a physical referent for the extra
information in $A$; the Aharonov--Bohm effect shows that the limit of
validity of this statement lies at the quantum level.
\end{remark}

\subsection{Complex (self-dual) formulation}

Because $\star\star=-1$ on 2-forms in Lorentzian signature, $\star$ has
eigenvalues $\pm\ii$ and $F$ splits into a self-dual and an anti-self-dual
part, which are complex conjugates of each other.  In three-dimensional
language this is the \emph{Riemann--Silberstein vector}
\begin{equation}\label{eq:RS}
\mathbf F=\vE+\ii\vB,\qquad
\divg\mathbf F=4\pi\rho,\qquad
\ii\,\pa_t\mathbf F=c\,\curl\mathbf F-4\pi\ii\,\vj .
\end{equation}
Two real vector equations have become one complex vector equation, of
Schr\"odinger form.  Under a rotation $\mathbf F$ transforms as a
vector; under a boost along the unit vector $\vn$ with rapidity
$\eta$ ($\tanh\eta=\beta$),
\begin{equation}
\mathbf F'_\parallel=\mathbf F_\parallel,\qquad
\mathbf F'_\perp=\cosh\eta\,\mathbf F_\perp-\ii\sinh\eta\;\vn\times\mathbf F_\perp ,
\end{equation}
which is a complex rotation about $\vn$ by the imaginary angle $\ii\eta$.
Thus $SO(3,1)^\uparrow_+\cong SO(3,\mathbb C)$ acts on $\mathbf F\in\mathbb C^3$:
the six real components of the field form \emph{one} complex vector of the
complexified rotation group.  The two invariants \eqref{eq:invariants}
are the real and imaginary parts of the only invariant of
$SO(3,\mathbb C)$,
\begin{equation}
\mathbf F\cdot\mathbf F=(\vE^2-\vB^2)+2\ii\,\vE\cdot\vB .
\end{equation}
In two-spinor language the same statement is that $F_{\mu\nu}$
corresponds to a symmetric spinor $\phi_{AB}$ (three complex components)
and its conjugate, in the representation $(1,0)\oplus(0,1)$ of
$SL(2,\mathbb C)$; the source-free equations are $\nabla^{AA'}\phi_{AB}=0$.

\subsection{Algebraic classification of fields}

Let $I_1=\vE^2-\vB^2$ and $I_2=\vE\cdot\vB$.  The matrix
$F^\mu{}_\nu$ has characteristic polynomial
$\det(\lambda-F^\mu{}_\nu)=\lambda^4-I_1\lambda^2-I_2^2$.  Consequently:
\begin{itemize}
\item \textbf{Non-null fields} ($I_1\ne0$ or $I_2\ne0$).  If $I_2\neq0$ there is a frame in which
  $\vE\parallel\vB$, with
  $E^2=\tfrac12\big(\sqrt{I_1^2+4I_2^2}+I_1\big)$ and
  $B^2=\tfrac12\big(\sqrt{I_1^2+4I_2^2}-I_1\big)$; the eigenvalues of
  $F^\mu{}_\nu$ are $\{\pm E,\ \pm\ii B\}$ (a boost eigenplane and a
  rotation eigenplane).  If $I_2=0$, $I_1>0$ there is a frame with
  $\vB=0$ (purely electric); if $I_2=0$, $I_1<0$ there is a frame with
  $\vE=0$ (purely magnetic).
\item \textbf{Null fields} ($I_1=I_2=0$): $\vE\perp\vB$ and $|\vE|=|\vB|$ in
  every frame; no frame exists in which the field vanishes or is purely
  electric or magnetic.  $F_{\mu\nu}=k_\mu a_\nu-k_\nu a_\mu$ for a null
  vector $k^\mu$ (the \emph{principal null direction}), which is the
  direction of $\vE\times\vB$.  Plane electromagnetic waves are null
  fields; a boost changes their frequency and amplitude but never their
  null character.
\end{itemize}
For a continuous solution the set of non-null points is open, so the
classification partitions spacetime into an open region where the field
is non-null and a closed set where it is null; a plane wave lies
entirely in the second.

\subsection{Symmetries}

\paragraph{Poincaré invariance.}  The equations are covariant under the ten-parameter
Poincaré group; Noether's theorem applied to \eqref{eq:action} yields
$P^\mu$ (four charges: $W/c,\ \mathbf P$) and $M^{\mu\nu}$ (six:
angular momentum $\mathbf J$ and the boost charges
$\mathbf K=\int(\vx\,u-t\,\vS)\,\dd^3x$).

\paragraph{Conformal invariance.}  Because $T^\mu{}_\mu=0$ (equivalently
because $\star$ on 2-forms in four dimensions depends only on the
conformal class of the metric), the free Maxwell equations are invariant under
the conformal group $SO(4,2)$ (fifteen parameters: the ten Poincaré
transformations, one dilatation, four special conformal transformations)
with $F_{\mu\nu}$ of conformal weight two.  This extends the conserved
charges to fifteen.  It is a four-dimensional accident: it fails
in any other dimension for the standard Maxwell action.

\paragraph{Duality.}  In the absence of sources,
$\vE\to\vE\cos\theta+\vB\sin\theta,\ \vB\to\vB\cos\theta-\vE\sin\theta$
maps solutions into solutions; in the complex notation \eqref{eq:RS}
it is $\mathbf F\to e^{-\ii\theta}\mathbf F$, a $U(1)$ which commutes
with the whole Poincaré group but does \emph{not} come from a
manifestly covariant local symmetry of the Lagrangian \eqref{eq:action}
(the potentials $A_\mu$ transform non-locally).  The conserved quantity
is the \emph{helicity}, the difference between the numbers of
positive- and negative-helicity photons, written in the mode amplitudes
of Section~\ref{sec:mx-free} as $\int\dd^3k\,(2\pi)^{-3}(|b_+|^2-|b_-|^2)$.  Electric sources break duality unless magnetic charges are added, so
that the pair of charges rotates together with $(\vE,\vB)$.  No magnetic
charge has been observed: the absence of monopoles, i.e.\
\eqref{eq:gauss-B}, is an \emph{empirical} statement, while the
equation $\dd F=0$ is the \emph{ideal-world} encoding of it.

% ==========================================================================
\section{Counting the degrees of freedom}\label{sec:mx-dof}
% ==========================================================================

\subsection{What is being counted}

``How many degrees of freedom does the electromagnetic field have?'' has
a single answer, \emph{two per point of space}, but the number is
reached by subtracting from larger numbers that are all natural to
quote ($6$ field components, $4$ potentials, $8$ equations), and the
subtractions are easily confused.  We distinguish three notions.
\begin{enumerate}
\item The number of \emph{components} of the unknown fields ($6$ for
  $(\vE,\vB)$, $4$ for $A^\mu$).
\item The number of \emph{free functions of the three space variables}
  in the general solution, i.e.\ the size of the freely specifiable
  Cauchy data.  For the Maxwell field this is $4$.
\item The number of \emph{degrees of freedom} $N$: half the previous
  number, when the dynamics is first order in time and the free data
  come in canonically conjugate pairs $(q,p)$.  For the Maxwell field
  $N=2$.
\end{enumerate}
The factor of two between (2) and (3) is the familiar one between the
number of initial data of a second-order equation (position and
velocity) and its number of degrees of freedom.  We now count in
several independent ways; they all give $N=2$.

\subsection{\texorpdfstring{Counting in terms of $\vE$ and $\vB$}{Counting in terms of E and B}}

There are $6$ unknown functions and $8$ equations \eqref{eq:maxwell-3d}.
Two relations reduce the number of independent equations:
\begin{align}
\divg\big[\text{\eqref{eq:faraday}}\big]-\frac1c\,\pa_t\big[\text{\eqref{eq:gauss-B}}\big]&\equiv0
&&\text{(an \emph{identity}: $\divg\curl\equiv0$),}\\
\divg\big[\text{\eqref{eq:ampere}}\big]+\frac1c\,\pa_t\big[\text{\eqref{eq:gauss-E}}\big]
&=-\frac{4\pi}{c}\big(\pa_t\rho+\divg\vj\big)
&&\text{(a \emph{compatibility condition} on the sources).}
\end{align}
Thus the $8$ equations contain $8-1-1=6$ independent conditions for $6$
unknowns: the system is determined, not overdetermined.  Concretely,
there are $6$ evolution equations \eqref{eq:evolution} and $2$
constraints.  The constraints are first integrals of the evolution
(Section~\ref{sec:mx-eqs}); hence the free Cauchy data are
\begin{equation}
\underbrace{6}_{\vE,\ \vB}-\underbrace{2}_{\text{constraints}}=4
\quad\text{functions on a spatial slice}=2\ \text{degrees of freedom}\times(q,p).
\end{equation}

\subsection{Fourier (characteristic) analysis}

Insert $\vE,\vB\propto e^{\ii(\vk\cdot\vx-\omega t)}$ into the source-free
evolution equations \eqref{eq:evolution}:
\begin{equation}\label{eq:eigen}
\omega\,\vE=-c\,\vk\times\vB,\qquad \omega\,\vB=c\,\vk\times\vE .
\end{equation}
Eliminating $\vB$ gives $\omega^2\vE=c^2\big[k^2\vE-\vk(\vk\cdot\vE)\big]$.
The $6\times6$ system therefore has the eigenfrequencies
\begin{equation}
\omega=+ck\ (\times2),\qquad \omega=-ck\ (\times2),\qquad \omega=0\ (\times2).
\end{equation}
For $\omega=\pm ck$ the eigenvectors are transverse, $\vE\perp\vk$,
$\vB=\pm\hat{\vk}\times\vE$: two polarizations for each sign of the
frequency.  The two eigenvectors with $\omega=0$ are the longitudinal
ones, $\vE\parallel\vk$, $\vB\parallel\vk$: they are precisely the
Fourier components of $\divg\vE$ and $\divg\vB$, i.e.\ the \emph{constraints},
and the vanishing frequency is the Fourier-space statement of
\eqref{eq:propagate-B}--\eqref{eq:propagate-E}: \emph{constraint
violations do not propagate; they are frozen}.  Of the six modes, four
propagate: $2$ polarizations $\times$ $(\pm\omega)$.  Reality of the
field relates the mode at $(\vk,\omega)$ to the one at
$(-\vk,-\omega)$, so per wave vector $\vk$ there remain two independent
complex amplitudes, i.e.\ four real numbers, i.e.\ $2$ degrees of
freedom.

\subsection{Counting with potentials}

The potentials have $4$ components, of which two are removed:
\begin{itemize}
\item \textbf{Lorenz gauge.}  The four wave equations \eqref{eq:lorenz-wave}
  are subject to $\pa_\mu A^\mu=0$ ($-1$) and have the residual gauge
  freedom $\Box\chi=0$ ($-1$): $4-1-1=2$.  For a plane wave
  $A^\mu=\epsilon^\mu e^{-\ii k\cdot x}$ with $k^2=0$ the Lorenz condition
  $k_\mu\epsilon^\mu=0$ admits three polarization vectors (including
  $\epsilon^\mu\propto k^\mu$, since $k$ is null), and the residual gauge
  transformation $\epsilon^\mu\to\epsilon^\mu+\alpha k^\mu$ removes one:
  $3-1=2$.  The ``timelike'' and ``longitudinal'' photons are the two
  that are not physical.
\item \textbf{Coulomb (radiation) gauge.}  $\divg\vA=0$.  Then $\varphi$ obeys
  the Poisson equation $\nabla^2\varphi=-4\pi\rho$ and is not an independent
  dynamical variable, while $\vA=\vA_\perp$ has $3-1=2$ independent
  components.
\end{itemize}

\subsection{Hamiltonian counting (Dirac--Bergmann)}

With $\mathcal L$ from \eqref{eq:action} the canonical momenta are
\begin{equation}
\Pi^i=\frac{\pa\mathcal L}{\pa\dot A^i}=-\frac{E^i}{4\pi c},\qquad
\Pi_\varphi=\frac{\pa\mathcal L}{\pa\dot\varphi}=0 .
\end{equation}
The vanishing of $\Pi_\varphi$ is a \emph{primary constraint}.  The
canonical Hamiltonian is
\begin{equation}\label{eq:H-can}
H=\int\dd^3x\Big[\frac{\vE^2+\vB^2}{8\pi}-\frac1c\,\vj\cdot\vA-\varphi\,\mathcal G\Big],\qquad
\mathcal G\equiv\frac{\divg\vE}{4\pi}-\rho ,
\end{equation}
after an integration by parts; $\varphi$ appears linearly, as a
\emph{Lagrange multiplier}.  Preservation of $\Pi_\varphi\approx0$,
$\dot\Pi_\varphi=-\delta H/\delta\varphi=\mathcal G\approx0$, gives the
\emph{secondary constraint}: Gauss's law.  Further, $\dot{\mathcal G}=0$
holds by virtue of \eqref{eq:continuity}, so the algorithm terminates.
All brackets among the constraints vanish, so both are
\emph{first class}.  With the phase space $(A^\mu,\Pi_\mu)$ of
dimension $8$ per point,
\begin{equation}
N=\frac12\big(8-2\times 2\big)=2 ,
\end{equation}
since each first-class constraint removes two phase-space dimensions:
one by restricting to the constraint surface, one by identifying the
points of a gauge orbit.  The first-class constraint $\mathcal G$ generates gauge
transformations: for
$G[\lambda]=c^{-1}\!\int\lambda\,\mathcal G\,\dd^3x$ and the canonical bracket
$\{A^i(\vx),E^j(\vy)\}=-4\pi c\,\delta^{ij}\delta^3(\vx-\vy)$,
$\delta_\lambda\vA=\{\vA,G[\lambda]\}=\grad\lambda$ (and
$\delta_\lambda\varphi=-c^{-1}\pa_t\lambda$ through $\Pi_\varphi$),
which is \eqref{eq:gauge} with $\chi=-\lambda$.  The physical phase space is the
quotient of the surface $\{\mathcal G=0\}$ by these transformations; its
symplectic form
\begin{equation}
\Omega=\int\dd^3x\ \delta\Pi^i\wedge\delta A^i
=-\frac1{4\pi c}\int\dd^3x\ \delta E^i\wedge\delta A^i
\end{equation}
is well defined on it because gauge directions are null vectors of $\Omega$
restricted to $\{\mathcal G=0\}$.  Decomposing
$\vE=\vE_\parallel+\vE_\perp$, $\vA=\vA_\parallel+\vA_\perp$ (Helmholtz),
$\mathcal G=0$ fixes $\vE_\parallel=-\grad\varphi_C$ with
$\nabla^2\varphi_C=-4\pi\rho$, and $\vA_\parallel$ is pure gauge, leaving
the canonical pairs $(\vA_\perp,\ -\vE_\perp/4\pi c)$.  The energy then
splits,
\begin{equation}\label{eq:H-split}
H=\int\dd^3x\,\frac{\vE_\perp^2+\vB^2}{8\pi}\ +\ \frac12\int\rho\,\varphi_C\,\dd^3x\ -\ \frac1c\int\vj\cdot\vA_\perp\,\dd^3x ,
\end{equation}
into radiation energy, Coulomb energy (instantaneous, determined by
the instantaneous charges), and interaction.  (Here the Coulomb gauge $\vA_\parallel=0$ was used; the term
$-\varphi\,\mathcal G$ vanishes on the constraint surface, and
$\int\vE_\parallel^2\,\dd^3x/8\pi=\tfrac12\int\rho\,\varphi_C\,\dd^3x$.)

\subsection{Group-theoretic counting and the massive comparison}

In the particle language, the quanta of the free field are massless
particles of spin $1$.  A massless one-particle state with momentum
$k^\mu=(\omega/c,0,0,k)$ is classified by the little group $ISO(2)$ of
the Euclidean plane; finite-component representations that are
physically acceptable have the translations acting trivially and are
labelled by the helicity $h$.  The photon has $h=\pm1$: \emph{two}
states, regardless of the four components of $A^\mu$.  A
Lorentz-covariant vector field cannot describe a massless helicity
$\pm1$ particle without a redundancy; the gauge freedom is exactly that
redundancy.

Contrast the massive (Proca) field, $\pa_\mu F^{\mu\nu}+\kappa^2A^\nu=(4\pi/c)j^\nu$ with
$\kappa=mc/\hbar$: there is no gauge freedom, the divergence gives
$\kappa^2\pa_\nu A^\nu=0$ (for conserved $j$) as a genuine \emph{constraint}, so
$4-1=3$ degrees of freedom, matching the little group $SO(3)$ of a massive
particle.  The gap between $3$ and $2$ is the longitudinal mode.  It decouples
from conserved currents as $m\to0$, so the massless limit of
observables is smooth, although the number of degrees of freedom jumps
from $3$ to $2$; its absence is equivalent to gauge invariance.

\subsection{Summary of the counts}

\begin{center}
\begin{tabular}{@{}lccc@{}}
\toprule
Method & Start & Subtract & Result\\
\midrule
Cauchy data for $(\vE,\vB)$ & $6$ functions & $2$ constraints & $4=2\times(q,p)$\\
Fourier / characteristics & $6$ eigenmodes & $2$ with $\omega=0$ (constraints) & $4=2\times(\pm\omega)$\\
Lorenz gauge & $4$ components & Lorenz condition, residual gauge & $2$\\
Coulomb gauge & $4$ components & $\varphi$ slaved, $\divg\vA=0$ & $2$\\
Dirac--Bergmann & $8$ phase-space dim. & $2$ first-class constraints ($2\times2$) & $4=2\times 2$\\
Representation theory & --- & little group $ISO(2)$, $h=\pm1$ & $2$\\
\bottomrule
\end{tabular}
\end{center}

% ==========================================================================
\section{Free solutions}\label{sec:mx-free}
% ==========================================================================

\subsection{Wave equations and their limitation}

Taking the curl of \eqref{eq:faraday} and \eqref{eq:ampere} and using
\eqref{eq:gauss-E}, \eqref{eq:gauss-B} one finds
\begin{equation}\label{eq:wave-EB}
\Box\vE=-4\pi\Big(\grad\rho+\frac1{c^2}\,\pa_t\vj\Big),\qquad
\Box\vB=\frac{4\pi}{c}\,\curl\vj ,
\end{equation}
so in vacuum each Cartesian component of $\vE,\vB$ obeys the wave
equation with speed $c$.  The converse is false, and the point is
important: the wave equations are \emph{consequences} of Maxwell's
equations, not equivalents.  A second-order equation for a vector field has
data $(\vE_0,\dot\vE_0)$: $6$ functions for $\vE$ and $6$ for $\vB$;
Maxwell's equations fix $\dot\vE_0$ and $\dot\vB_0$ in terms of the curls
of the other field and impose two divergence conditions, cutting the
twelve data to four.

\subsection{Plane waves}

The simplest free solutions are
$\vE=\operatorname{Re}\,[\mathbf e\,e^{\ii(\vk\cdot\vx-\omega t)}]$ with
\begin{equation}\label{eq:plane}
\omega=c|\vk|,\qquad \vk\cdot\mathbf e=0,\qquad
\vB=\hat{\vk}\times\vE,\qquad \vS=c\,u\,\hat{\vk}.
\end{equation}
The field is transverse and null ($I_1=I_2=0$), the energy flows at speed $c$
along $\hat{\vk}$, and $|\vS|=cu$ saturates the bound \eqref{eq:S-bound}.
Let $\mathbf e_1,\mathbf e_2,\hat{\vk}$ be a right-handed orthonormal
triple.  The circular polarization (\emph{helicity}) vectors
\begin{equation}
\mathbf e_\pm=\frac{\mathbf e_1\pm\ii\,\mathbf e_2}{\sqrt2},\qquad
\ii\,\vk\times\mathbf e_\pm=\pm k\,\mathbf e_\pm ,\qquad
\mathbf e_\lambda^*\cdot\mathbf e_{\lambda'}=\delta_{\lambda\lambda'},
\end{equation}
diagonalize the curl and the generator of rotations about $\hat{\vk}$;
with the convention $e^{\ii(\vk\cdot\vx-\omega t)}$ the vector
$\mathbf e_+$ rotates counter-clockwise about $\hat{\vk}$, i.e.\ has
positive helicity.  A general polarization (ellipse) is a complex
combination of $\mathbf e_+$ and $\mathbf e_-$.

\subsection{General solution in the Coulomb gauge}

In the radiation gauge with $\rho=0,\vj=0$ one has $\varphi=0$,
$\divg\vA=0$, $\Box\vA=0$, and the general real solution is
\begin{equation}\label{eq:modes}
\vA(t,\vx)=\int\!\frac{\dd^3k}{(2\pi)^3}\sum_{\lambda=\pm}\sqrt{\frac{2\pi c}{k}}\Big[b_\lambda(\vk)\,\mathbf e_\lambda(\vk)\,e^{\ii(\vk\cdot\vx-\omega t)}+\text{c.c.}\Big],
\end{equation}
which gives
\begin{align}
\vE&=\int\!\frac{\dd^3k}{(2\pi)^3}\sum_\lambda\ii\sqrt{2\pi ck}\Big[b_\lambda\,\mathbf e_\lambda e^{\ii(\vk\cdot\vx-\omega t)}-\text{c.c.}\Big],\\
\vB&=\int\!\frac{\dd^3k}{(2\pi)^3}\sum_\lambda\lambda\sqrt{2\pi ck}\Big[b_\lambda\,\mathbf e_\lambda e^{\ii(\vk\cdot\vx-\omega t)}+\text{c.c.}\Big].
\end{align}
All three $\vk$ components range over $\mathbb R^3$: the negative-frequency
term at $\vk$ is the complex conjugate of the positive-frequency term at
$-\vk$, so there is no double counting.  The two complex functions
$b_\pm(\vk)$ of three variables, i.e.\ four real functions, are the free data.  Substitution into
\eqref{eq:densities} and the orthogonality of the $e^{\ii\vk\cdot\vx}$
(the terms $\propto b\,b$ cancel between $\vE^2$ and $\vB^2$) gives the
conserved quantities in diagonal form:
\begin{equation}\label{eq:W-P-S}
W=\int\!\frac{\dd^3k}{(2\pi)^3}\sum_\lambda\omega\,|b_\lambda|^2,\qquad
\mathbf P=\int\!\frac{\dd^3k}{(2\pi)^3}\sum_\lambda\vk\,|b_\lambda|^2,\qquad
\mathbf J_{\rm spin}=\int\!\frac{\dd^3k}{(2\pi)^3}\sum_\lambda\lambda\,\hat{\vk}\,|b_\lambda|^2 .
\end{equation}
Since $\dot b_\lambda=-\ii\omega\,b_\lambda$ follows from $H=W$, the
canonical bracket is $\{b_\lambda(\vk),b^*_{\lambda'}(\vk')\}=-\ii(2\pi)^3\delta_{\lambda\lambda'}\delta^3(\vk-\vk')$.
Each pair $(\vk,\lambda)$ is an independent harmonic oscillator, with
\emph{action variable} $|b_\lambda|^2$ and \emph{angle variable}
$\arg b_\lambda$; the free field is an integrable system with an
infinite set of conserved action variables, the energy, momentum and
helicity being linear combinations of them.  For each mode $\vk|b|^2=(\omega|b|^2/c)\,\hat{\vk}$: momentum equals
energy$/c$ along $\hat{\vk}$, the mode-by-mode form of $\vS=cu\hat{\vk}$.

\subsection{Spherical waves}

For solutions with definite angular momentum $(\ell,m)$, $\ell\ge1$, a
complete set is given by two families built on vector spherical
harmonics: \emph{electric} (TM, $B_r=0$) and \emph{magnetic} (TE, $E_r=0$),
distinguished by parity $(-1)^{\ell}$ and $(-1)^{\ell+1}$.  There is no
$\ell=0$ radiating solution: a spherically symmetric configuration
produces no radiation (Gauss's law and symmetry force $\vB=0$ and the
Coulomb field), an analogue of Birkhoff's theorem in
electrodynamics; the reason there are exactly two independent families
for each $(\ell,m)$ is again the two helicities.

% ==========================================================================
\section{The initial-value (Cauchy) problem}\label{sec:mx-cauchy}
% ==========================================================================

\subsection{Statement of the problem}

\begin{quote}
\textbf{Problem.}  Given $\rho(t,\vx)$ and $\vj(t,\vx)$ satisfying
\eqref{eq:continuity} and fields $\vE_0(\vx),\vB_0(\vx)$ at $t=0$
satisfying $\divg\vE_0=4\pi\rho(0,\cdot)$ and $\divg\vB_0=0$, find $\vE,\vB$ for all
$t$ solving \eqref{eq:maxwell-3d}.
\end{quote}
The data are specified on the spacelike hyperplane $t=0$, the evolution
equations \eqref{eq:evolution} are first order in $t$, and the
constraints are preserved by \eqref{eq:propagate-B}--\eqref{eq:propagate-E}.  It remains to decide whether the
problem is well posed (in the sense of Hadamard: existence, uniqueness,
continuous dependence on the data).

\subsection{Symmetric hyperbolic form and well-posedness}

Write $u=(\vE,\vB)^{\mathsf T}\in\mathbb R^6$.  Then \eqref{eq:evolution} is
\begin{equation}\label{eq:symhyp}
\pa_t u=\sum_{b=1}^{3}A^b\,\pa_b u-4\pi\,(\vj,\mathbf 0)^{\mathsf T},\qquad
(A^b)_{E_aB_c}=(A^b)_{B_cE_a}=c\,\varepsilon_{abc},
\end{equation}
all other entries of $A^b$ being zero.  (Indeed
$\pa_tE_a=c\,\varepsilon_{abc}\pa_bB_c$ and
$\pa_tB_c=-c\,\varepsilon_{cba}\pa_bE_a=c\,\varepsilon_{abc}\pa_bE_a$.)
Each $A^b$ is a \emph{real symmetric} constant matrix.  A first-order system with
real symmetric constant coefficient matrices is \emph{symmetric
hyperbolic} (Friedrichs).  Its principal symbol for a covector
$\xi$ is $\sum_bA^b\xi_b$, with eigenvalues $0,0,\pm c|\xi|,\pm c|\xi|$:
the characteristic speeds are $0$ and $\pm c$, the zero speed again being the
constraint sector, and the maximal speed is $c$.  The operator
$\mathcal A=\sum A^b\pa_b$ is skew-adjoint on
$L^2(\mathbb R^3;\mathbb R^6)$,
\begin{equation}
\langle u,\mathcal Au\rangle=c\int(\vE\cdot\curl\vB-\vB\cdot\curl\vE)\,\dd^3x
=c\int\divg(\vB\times\vE)\,\dd^3x=0,
\end{equation}
so $\mathcal A$ generates a unitary group (Stone's theorem) and the energy
$W=\|u\|^2_{L^2}/8\pi$ is conserved for $\vj=0$.  Duhamel's formula
handles the source.

\begin{theorem}[Well-posedness in $H^s$]\label{thm:wellposed}
Let $s\ge0$, $(\vE_0,\vB_0)\in H^s(\mathbb R^3)^6$ and
$\vj\in L^1_{\rm loc}(\mathbb R;H^s(\mathbb R^3)^3)$.  Then there is a unique
solution $u\in C(\mathbb R;H^s)$ of \eqref{eq:symhyp} with
$u(0)=(\vE_0,\vB_0)$, depending continuously on the data, and
\begin{equation}\label{eq:energy-estimate}
\|u(t)\|_{H^s}\le\|u(0)\|_{H^s}+4\pi\Big|\int_0^t\|\vj(\tau)\|_{H^s}\,\dd\tau\Big| .
\end{equation}
If, in addition, a charge density $\rho$ with $\pa_t\rho+\divg\vj=0$ is
given and the constraints \eqref{eq:gauss-E}, \eqref{eq:gauss-B} hold at
$t=0$, then they hold for all $t$.
\end{theorem}
\noindent\emph{Proof sketch.}  Differentiating $\tfrac12\|u\|^2$ and using
skew-adjointness of $\mathcal A$,
$\tfrac12\,\dd_t\|u\|^2=-4\pi\int\vj\cdot\vE\,\dd^3x\le4\pi\|\vj\|\,\|u\|$,
so $\dd_t\|u\|\le4\pi\|\vj\|$, which integrates to
\eqref{eq:energy-estimate} for $s=0$.  Because the coefficients are
constant, $\mathcal A$ commutes with $\pa^\alpha$, so the same estimate
holds for every Sobolev norm.  Existence follows from Stone's theorem
and Duhamel, $u(t)=e^{t\mathcal A}u_0-4\pi\int_0^te^{(t-\tau)\mathcal A}(\vj(\tau),\mathbf0)\,\dd\tau$,
uniqueness from the estimate applied to the difference of two
solutions.  Propagation of constraints is \eqref{eq:propagate-B}--\eqref{eq:propagate-E}.
$\square$

The estimate expresses stability: small changes of the data and of
the sources produce small changes of the solution, with no loss of
derivatives, and with a growth that is at most linear in time for
bounded sources.

\subsection{Finite propagation speed and domain of dependence}

The $L^2$ theory proves more: information propagates at speed at most $c$.
Let $B(\vx_0,R-ct)$ be the ball of radius $R-ct$ around $\vx_0$, and
$W_R(t)=\int_{B(\vx_0,R-ct)}u\,\dd^3x$ its energy.  Then, using
\eqref{eq:poynting} and the fact that the boundary recedes at speed $c$,
\begin{equation}
\frac{\dd W_R}{\dd t}
=-\oint_{\pa B}\big(\vS\cdot\vn+c\,u\big)\,\dd a-\int_{B}\vj\cdot\vE\,\dd^3x
\le-\int_B\vj\cdot\vE\,\dd^3x ,
\end{equation}
since $\vS\cdot\vn+cu\ge0$ by \eqref{eq:S-bound}.

\begin{proposition}[Domain of dependence]\label{prop:local-energy}
If $\vj=0$ in the cone $\mathcal D=\{(t,\vx):|\vx-\vx_0|<R-c|t|\}$ and
$\vE_0=\vB_0=0$ on $B(\vx_0,R)$, then $\vE=\vB=0$ throughout $\mathcal D$.
Equivalently, the fields at $(t,\vx)$ depend only on the data in the
ball $B(\vx,c|t|)$ and on the sources inside the backward light cone of
$(t,\vx)$.
\end{proposition}
\noindent\emph{Proof.}  $W_R(0)=0$, $W_R\ge0$ and $\dd W_R/\dd t\le0$, hence
$W_R\equiv0$ for $0\le t<R/c$; the case $t<0$ follows from time
reversal (last subsection of this section).  $\square$

\begin{figure}[t]
\centering
\begin{tikzpicture}[>=stealth,scale=1.0]
\draw[->] (-4,0) -- (4.3,0) node[right] {$\vx$};
\draw[->] (0,-0.1) -- (0,3.6) node[above] {$ct$};
\fill[gray!20] (-3,0) -- (3,0) -- (0,3) -- cycle;
\draw[very thick] (-3,0) -- (3,0);
\draw[thick] (-3,0) -- (0,3) -- (3,0);
\node at (0,0.95) {$\mathcal D$};
\node[below] at (-3,0) {$\vx_0-R$};
\node[below] at (3,0) {$\vx_0+R$};
\node[below left] at (0,0) {$\vx_0$};
\fill (0,3) circle (1.5pt);
\node[above right] at (0.05,3) {$ct=R$};
\node[right] at (3.15,0.55) {data on $t=0$};
\draw[->,thin] (3.15,0.55) -- (2.6,0.07);
\end{tikzpicture}
\caption{Domain of dependence $\mathcal D$ of the ball $B(\vx_0,R)$: if the data vanish on the
thick segment (and there are no sources in $\mathcal D$), the solution vanishes in the shaded
cone, whose sides are light rays.}\label{fig:domain}
\end{figure}

Proposition~\ref{prop:local-energy} is the PDE statement of
\emph{relativistic causality} for the free field: the light cone is the
boundary of the domain of dependence, and $c$ is the universal maximal speed
at which $\vE,\vB$ can carry information.

\subsection{Explicit solution and Huygens' principle}

For $\vj=0$ each Cartesian component of $\vE,\vB$ solves the wave
equation $u_{tt}=c^2\nabla^2u$ with data $(u_0,u_1)$ obtained from
the evolution equations: $\dot\vE_0=c\curl\vB_0$, $\dot\vB_0=-c\curl\vE_0$.
Kirchhoff's formula gives, for $t>0$,
\begin{equation}\label{eq:kirchhoff}
\begin{aligned}
\vE(t,\vx)&=\pa_t\big[t\,M_{ct}\vE_0\big](\vx)+c\,t\,M_{ct}\big[\curl\vB_0\big](\vx),\\
\vB(t,\vx)&=\pa_t\big[t\,M_{ct}\vB_0\big](\vx)-c\,t\,M_{ct}\big[\curl\vE_0\big](\vx),
\end{aligned}
\qquad
M_rf(\vx)=\frac1{4\pi}\oint_{|\boldsymbol\omega|=1}f(\vx+r\boldsymbol\omega)\,\dd\Omega .
\end{equation}
Thus the fields at $(t,\vx)$ depend \emph{only on the data on the sphere
$|\vx'-\vx|=ct$}, not on the ball inside it: this is the \emph{strong
Huygens principle}.  A localized pulse is received as a localized pulse,
with a sharp front and a sharp rear, and no tail.  In three space
dimensions this holds because the dimension is odd and larger than one;
it fails in two (e.g.\ a line source) and for the Maxwell
equations in a medium with dispersion.  Huygens' principle is not a
consequence of causality alone: it is the special property that makes
the exact reconstruction of signals and the sharpness of radiation
fronts possible in free space.

\subsection{Solving the constraints on the initial slice}

The data on $t=0$ are not freely specifiable: they must satisfy
$\divg\vE_0=4\pi\rho_0$ and $\divg\vB_0=0$.  For fields in $L^2(\mathbb R^3)$
the orthogonal (Helmholtz) decomposition $\vF=\vF_\parallel+\vF_\perp$ is
given in Fourier space by $\widehat{\vF}_\parallel=\hat{\vk}(\hat{\vk}\cdot\widehat{\vF})$, $\widehat{\vF}_\perp=\widehat{\vF}-\widehat{\vF}_\parallel$.  Then
\begin{equation}
\vB_{0\parallel}=0,\qquad
\vE_{0\parallel}=-\grad\varphi_0,\quad \varphi_0(\vx)=\int\frac{\rho_0(\vx')}{|\vx-\vx'|}\,\dd^3x',
\end{equation}
and the free data are $\vE_{0\perp}$ and $\vB_0$: two functions each
(a divergence-free vector field on $\mathbb R^3$ has two independent
components), i.e.\ four, as in Section~\ref{sec:mx-dof}.  The longitudinal
(Coulomb) part of the field is slaved to the instantaneous charge
distribution; all the independent dynamics lives in the transverse
fields.  On a domain with non-trivial topology there are, in addition,
finite-dimensional spaces of harmonic fields (Section~\ref{sec:mx-bvp}).

\subsection{Cauchy problem for the potentials}

For the potentials, count in the Lorenz gauge.  The second-order system
\eqref{eq:lorenz-wave} has data $(A^\mu,\pa_tA^\mu)|_{t=0}$: $8$
functions.  The Lorenz condition is preserved by the evolution only if it
holds initially together with its time derivative, because
$\Box(\pa_\mu A^\mu)=(4\pi/c)\pa_\mu j^\mu=0$.  Computing the time derivative
with the $\nu=0$ equation $\Box\varphi=4\pi\rho$ gives
\begin{equation}
\frac1c\pa_t(\pa_\mu A^\mu)=4\pi\rho-\divg\vE ,
\end{equation}
so the Lorenz condition and its derivative at $t=0$ are equivalent to the
Lorenz condition and \emph{Gauss's law}: two conditions ($8\to6$).  The
residual gauge $\Box\chi=0$ with data $(\chi,\pa_t\chi)$ removes two more
functions ($6\to4$).  Again $4=2\times(q,p)$.  In the Coulomb gauge the
data are $(\vA_\perp,\pa_t\vA_\perp)$ (four functions), $\varphi_C$ is
computed from $\rho$ at the same instant, and $\vA_\perp$ obeys
\begin{equation}
\Box\vA_\perp=\frac{4\pi}{c}\,\vj_\perp,\qquad \vj_\perp=\vj-\frac1{4\pi}\grad\pa_t\varphi_C .
\end{equation}

\begin{remark}[The Coulomb gauge is instantaneous, the field is not]
In the Coulomb gauge $\varphi_C(t,\vx)=\int\rho(t,\vx')/|\vx-\vx'|\,\dd^3x'$ responds
\emph{instantaneously} to the charge distribution, and the transverse part
$\vA_\perp$ is nonlocal because it is sourced by $\vj_\perp$.  Yet
the gauge-invariant fields $\vE=-\grad\varphi_C-c^{-1}\pa_t\vA_\perp$ and
$\vB=\curl\vA_\perp$ are causal: the superluminal pieces of $\grad\varphi_C$ and of
$c^{-1}\pa_t\vA_\perp$ cancel in the sum.  The instantaneous action at
a distance belongs to the \emph{ideal world} of a particular
parametrization (the gauge), not to the real world of measurable
fields.  Causality is a property of the gauge-invariant quantities.
\end{remark}

\subsection{Time reversal and the arrow of time}

Under $t\to-t$, $\vE\to\vE$, $\vB\to-\vB$, $\rho\to\rho$, $\vj\to-\vj$ the equations
\eqref{eq:maxwell-3d} are invariant.  The Cauchy problem is therefore
well posed backwards as well, \eqref{eq:energy-estimate} holds for $t<0$, and
the free evolution is a unitary group, not merely a semigroup:
\emph{nothing in the field equations at fixed sources distinguishes
the future from the past}.  The distinction arises only when one decides
how a free field and the sources are related: the retarded and the
advanced solutions of Section~\ref{sec:mx-sources} are equally good solutions.

% ==========================================================================
\section{The source problem: Green functions and radiation}\label{sec:mx-sources}
% ==========================================================================

\subsection{Green functions of the wave operator}

In the Lorenz gauge the sourced problem reduces to $\Box A^\mu=(4\pi/c)j^\mu$.
Let $G$ be a fundamental solution, $\Box G=\delta^4(x)\equiv\delta(x^0)\delta^3(\vx)$.
Two of them are of special interest, the retarded and the advanced one,
\begin{equation}\label{eq:green}
G_{\rm ret}(x)=\frac{\delta(x^0-r)}{4\pi r}=\frac1{2\pi}\,\theta(x^0)\,\delta(x^2),\qquad
G_{\rm adv}(x)=\frac{\delta(x^0+r)}{4\pi r}=\frac1{2\pi}\,\theta(-x^0)\,\delta(x^2),
\end{equation}
$r=|\vx|$, $x^2=(x^0)^2-r^2$.  They are the Fourier integrals
$G(x)=\int\!\frac{\dd^4k}{(2\pi)^4}\,\frac{e^{-\ii k\cdot x}}{|\vk|^2-k_0^2}$ with
the prescription $k_0\to k_0\pm\ii0$ (poles displaced to the lower, respectively
upper, half plane).  Their support lies on the light cone: this is
Huygens' principle for the Green function.  Their difference
\begin{equation}
D=G_{\rm ret}-G_{\rm adv}=\frac1{2\pi}\,\operatorname{sgn}(x^0)\,\delta(x^2)
=\frac{\delta(x^0-r)-\delta(x^0+r)}{4\pi r}
\end{equation}
is a solution of the \emph{homogeneous} equation with
$D|_{x^0=0}=0$, $\pa_0D|_{x^0=0}=\delta^3(\vx)$, so the solution of the free equation
$\Box u=0$ with data $(u_0,\pa_0u_0)$ on $x^0=0$ is
$u(x^0,\vx)=\int\dd^3x'\,\big[D(x^0,\vx-\vx')\,\pa_0u_0(\vx')+\pa_0D(x^0,\vx-\vx')\,u_0(\vx')\big]$,
which is Kirchhoff's formula \eqref{eq:kirchhoff} written with
the Pauli--Jordan function $D$.

\subsection{The retarded solution}

The retarded potential of a conserved source is
\begin{equation}\label{eq:retarded}
\begin{aligned}
A^\mu_{\rm ret}(x)&=\frac{4\pi}{c}\int G_{\rm ret}(x-x')\,j^\mu(x')\,\dd^4x',\quad\text{i.e.}\\
\varphi&=\int\!\frac{\rho(t-|\vx-\vx'|/c,\vx')}{|\vx-\vx'|}\,\dd^3x',\qquad
\vA=\frac1c\int\!\frac{\vj(t-|\vx-\vx'|/c,\vx')}{|\vx-\vx'|}\,\dd^3x' .
\end{aligned}
\end{equation}
It satisfies the Lorenz condition because $\pa_\mu j^\mu=0$, and it depends
on the source only on the backward light cone of the field point
(sharply, with no tail).  A particular solution is not the general one:
the general solution of the inhomogeneous Maxwell equations is
\begin{equation}\label{eq:in-out}
A^\mu=A^\mu_{\rm ret}+A^\mu_{\rm in}=A^\mu_{\rm adv}+A^\mu_{\rm out},\qquad
\Box A^\mu_{\rm in}=\Box A^\mu_{\rm out}=0,
\end{equation}
where $A_{\rm in}$ is the incoming field (the field present before the sources
act) and $A_{\rm out}$ the outgoing one.  Their difference
\begin{equation}
A_{\rm out}-A_{\rm in}=A_{\rm ret}-A_{\rm adv}=\frac{4\pi}{c}\int D(x-x')\,j(x')\,\dd^4x'
\end{equation}
is a free field: the \emph{radiation field} of the source.  Retarded solution
$=$ Cauchy problem with zero data in the infinitely remote past;
advanced solution $=$ zero data in the infinitely remote future.  Both
are legitimate, the equations are time-reversal invariant, and the
asymmetry is an additional condition, on the data at $t\to\mp\infty$.

\begin{remark}[Real world vs.\ ideal world: free fields and remote sources]
A field with no sources \emph{anywhere} in spacetime is an object of the
ideal world.  In the real world a ``free'' field is a field whose sources
are far from the region and the time of interest: starlight is a free
field in the laboratory, produced by sources that are outside the
region.  The decomposition \eqref{eq:in-out} makes this explicit:
$A_{\rm in}$ is the retarded field of the sources of the
distant past, written as if they did not exist.  The Cauchy problem for
free fields and the source problem are therefore two descriptions of
the same thing, differing in where the line is drawn between ``the
system'' and ``the rest of the universe''; this is why the degrees of
freedom counted in Section~\ref{sec:mx-dof} are carried by free fields
although, in the real world, every free field originates in some
source.  Choosing the retarded solution is the assumption that the
region contains no incoming radiation from the future, an assumption
about the whole of spacetime that no local experiment can check.
\end{remark}

\subsection{The field of a point charge; radiation}

For a point charge $e$ on a prescribed worldline $\vx=\mathbf z(t)$ the
integral \eqref{eq:retarded} gives the Liénard--Wiechert potentials
\begin{equation}
\varphi=\frac{e}{\big[(1-\vn\cdot\vbeta)R\big]_{\rm ret}},\qquad
\vA=\frac{e\,\vbeta}{\big[(1-\vn\cdot\vbeta)R\big]_{\rm ret}},\qquad
\vbeta=\frac{\dot{\mathbf z}}{c},\ \ R=|\vx-\mathbf z|,\ \ \vn=\frac{\vx-\mathbf z}{R},
\end{equation}
all bracketed quantities evaluated at the retarded time $t_{\rm r}$ defined by
$t_{\rm r}+R(t_{\rm r})/c=t$.  The fields are
\begin{equation}\label{eq:LW}
\vE=e\Big[\frac{\vn-\vbeta}{\gamma^2(1-\vn\cdot\vbeta)^3R^2}\Big]_{\rm ret}
+\frac ec\Big[\frac{\vn\times\big((\vn-\vbeta)\times\dot\vbeta\big)}{(1-\vn\cdot\vbeta)^3R}\Big]_{\rm ret},\qquad
\vB=[\vn\times\vE]_{\rm ret}.
\end{equation}
The first term (``velocity field'') falls as $R^{-2}$ and exists for
uniform motion; the second (``acceleration'' or \emph{radiation}
field) falls as $R^{-1}$, is proportional to $\dot\vbeta$, and for large $R$
satisfies $\vE\perp\vn$, $\vB=\vn\times\vE$, $|\vE|=|\vB|$: it is
asymptotically a \emph{null field} in the sense of
Section~\ref{sec:mx-geom}.  The radiated power of a slowly moving charge is Larmor's,
and in general Liénard's:
\begin{equation}
P=\frac{2e^2}{3c^3}\,|\dot{\vv}|^2,\qquad
P=\frac{2e^2}{3c}\,\gamma^6\Big[\dot\vbeta^2-(\vbeta\times\dot\vbeta)^2\Big].
\end{equation}
For a small source of dipole moment $\mathbf d(t)$ (with wavelength much larger than the
size of the source) the far field is $\vB=\ddot{\mathbf d}(t-r/c)\times\vn/(c^2r)$,
$\vE=\vB\times\vn$, giving $\dd P/\dd\Omega=|\ddot{\mathbf d}|^2\sin^2\theta/(4\pi c^3)$
and $P=2|\ddot{\mathbf d}|^2/(3c^3)$.

\begin{remark}[Real world vs.\ ideal world: the point source]
A $\delta$-function charge is the ideal-world limit of a small charged
body; the field \eqref{eq:LW} diverges at the charge, the electrostatic
self-energy $\tfrac12\!\int\!\rho\varphi$ is infinite, and the action of
the field on its own source (the radiation reaction) requires either a
regularization of the source or the renormalization of the mass.  These
problems belong to the dynamics of the charge-plus-field system, not to
the equations of the field with given sources, and for that reason they are
postponed to the treatment of the relativistic charged particle.  The fields \eqref{eq:LW} are
nonetheless solutions of the Maxwell equations with the $\delta$-source
in the sense of distributions, since they are locally integrable; it is
the field energy density, which behaves as $R^{-4}$, that is not
integrable at the charge.
\end{remark}

% ==========================================================================
\section{Boundary and initial--boundary value problems}\label{sec:mx-bvp}
% ==========================================================================

\subsection{Interfaces and the perfect conductor}

At a surface $\mathcal S$ with unit normal $\vn$ carrying charge density
$\sigma$ and current $\vK$ the jump conditions \eqref{eq:jump}
hold.  They connect the fields on the two sides and determine the surface
sources from the fields, or vice versa.  The most important idealization
is the \emph{perfect conductor}: the field vanishes inside, and for the
exterior vacuum region, with $\vn$ pointing out of the conductor,
\begin{equation}\label{eq:PEC}
\vn\times\vE=0,\qquad \vn\cdot\vB=0,\qquad
\sigma=\frac{\vn\cdot\vE}{4\pi},\qquad \vK=\frac c{4\pi}\,\vn\times\vB .
\end{equation}
The first two are the boundary conditions of the vacuum problem; the last
two merely determine the induced charge and current, which are
\emph{outputs}.

\begin{remark}[Real world vs.\ ideal world: the perfect conductor]
No real metal is a perfect conductor.  Condition \eqref{eq:PEC} is the
limit $\sigma_{\rm c}\to\infty$ of a conductor with Gaussian conductivity $\sigma_{\rm c}$: the field
penetrates a skin depth $\delta=c/\sqrt{2\pi\sigma_{\rm c}\omega}$ and the
perfect-conductor condition is the limit $\delta\to0$ in which this
boundary layer is replaced by a surface.  It is a \emph{singular limit}:
the dissipation, which is present for every finite $\sigma_{\rm c}$ and gives the
cavity modes of the next subsections a finite quality factor, disappears in the limit, so
that the energy of the ideal cavity is exactly conserved.  A theorem for the ideal problem
is a statement about the model, and its relevance to a copper
cavity is valid only on time scales short compared with the
dissipation time.
\end{remark}

\subsection{The initial--boundary value problem in a bounded domain}

Let $\Omega\subset\mathbb R^3$ be a bounded domain with smooth boundary
$\pa\Omega$ and outward normal $\vn$.

\begin{quote}
\textbf{Problem.}  Find $\vE,\vB$ in $\Omega\times\mathbb R$ solving
\eqref{eq:maxwell-3d}, with $\vn\times\vE=0$ on $\pa\Omega$ for all $t$,
and with given initial data $\vE_0,\vB_0$ satisfying the constraints in $\Omega$
and the compatibility conditions $\vn\times\vE_0=0$, $\vn\cdot\vB_0=0$ on $\pa\Omega$.
\end{quote}
Three observations explain why the problem is well posed.
\begin{enumerate}
\item \emph{Energy balance.}  With $\vj=0$ and the Poynting theorem,
  \begin{equation}
  \frac{\dd}{\dd t}\int_\Omega u\,\dd^3x=-\oint_{\pa\Omega}\vS\cdot\vn\,\dd a
  =-\frac c{4\pi}\oint_{\pa\Omega}(\vn\times\vE)\cdot\vB\,\dd a=0 ,
  \end{equation}
  since $\vn\cdot(\vE\times\vB)=(\vn\times\vE)\cdot\vB$.  The boundary condition
  makes the energy a conserved quantity: the problem is
  \emph{conservative}, not merely dissipative.  The uniqueness proof is the
  same as in Section~\ref{sec:mx-cauchy}; in terms of
  operators, the Maxwell operator with the domain
  $\{\vE\in H_0(\curl),\ \vB\in H(\curl)\}$ is skew-adjoint on
  $L^2(\Omega)^6$, so it generates a unitary group.
\item \emph{Compatibility of the boundary condition with the evolution.}
  $\pa_t(\vn\cdot\vB)=-c\,\vn\cdot\curl\vE$ involves only the
  tangential trace $\vn\times\vE$ of $\vE$ on the boundary, so
  $\vn\times\vE=0$ for all $t$ preserves $\vn\cdot\vB=0$.  For smooth
  solutions the initial data must also satisfy the higher-order
  compatibility conditions, e.g.\ $\vn\times(\curl\vB_0-4\pi\vj_0/c)=0$
  (because $\vn\times\pa_t\vE=0$), and so on.
\item \emph{Finite propagation speed with reflection.}
  Proposition~\ref{prop:local-energy} remains true as long as the cone
  does not meet $\pa\Omega$; waves arriving at the wall are reflected and
  the solution is determined by the data and by the geometry.
\end{enumerate}

\subsection{Topology and the static sector}

Zero-frequency solutions play the role of the constraint modes of
Section~\ref{sec:mx-dof}, and in a bounded domain they are governed by
topology.  Static fields with $\vj=0$, $\rho=0$ in $\Omega$ obey
$\curl\vE=0$, $\divg\vE=0$, $\vn\times\vE=0$, and $\curl\vB=0$,
$\divg\vB=0$, $\vn\cdot\vB=0$.
\begin{itemize}
\item $\vE=\grad\phi$ with $\phi$ constant on each connected component of $\pa\Omega$.
  If $\pa\Omega$ has $m$ components there are $m-1$ independent such fields:
  the charges on the separate conductors (a conductor surrounded by
  another, as in a coaxial line) are separately conserved.
\item $\vB$ harmonic with $\vn\cdot\vB=0$: there are $b_1(\Omega)$ independent
  ones, where $b_1$ is the first Betti number of $\Omega$ (the number of independent
  non-contractible loops, e.g.\ $1$ for a toroidal cavity): the flux
  through each ``handle'' is conserved.
\end{itemize}
These finite-dimensional spaces are exactly the frozen modes that
carry the topological part of the constraint sector; on all of
$\mathbb R^3$ they are absent.

\subsection{Cavity modes}

For $\vj=0$ separate $\vE=\operatorname{Re}[\vE_n(\vx)e^{-\ii\omega_nt}]$.
Then
\begin{equation}\label{eq:cavity}
\curl\curl\vE_n=k_n^2\vE_n,\qquad\divg\vE_n=0,\qquad\vn\times\vE_n=0\ \text{on }\pa\Omega,\qquad
\omega_n=ck_n,\ \ \vB_n=\frac1{k_n}\curl\vE_n .
\end{equation}
The operator $\curl\curl$ on divergence-free fields with this boundary condition has a compact
resolvent, hence a discrete spectrum $0<k_1\le k_2\le\dots\to\infty$ and
a complete orthonormal family $\{\vE_n\}$.  Expanding
$\vE=\sum e_n(t)\vE_n$, $\vB=\sum b_n(t)\vB_n$ (plus the static fields above)
the Maxwell equations become
\begin{equation}
\dot e_n=ck_n\,b_n-4\pi j_n,\qquad \dot b_n=-ck_n\,e_n,\qquad
\ddot e_n+\omega_n^2e_n=-4\pi\dot j_n,\qquad j_n=\int_\Omega\vj\cdot\vE_n\,\dd^3x,
\end{equation}
with energy $\sum(e_n^2+b_n^2)/8\pi$: the field in a cavity is a countable
set of independent, driven harmonic oscillators.  The degrees of
freedom, $2$ per point of the continuum, have become a discrete set of
oscillators; the number with frequency below $\omega$ is given for large $\omega$
by Weyl's law
\begin{equation}
N(\omega)\simeq 2\times\frac{V}{(2\pi)^3}\cdot\frac{4\pi}3\Big(\frac\omega c\Big)^3
=\frac{V\omega^3}{3\pi^2c^3},
\end{equation}
two polarizations per cell of volume $(2\pi)^3/V$ in $\vk$-space, which is
the count of Section~\ref{sec:mx-dof} in disguise.

\begin{example}[Rectangular cavity]
For a box $[0,a]\times[0,b]\times[0,d]$ with perfectly conducting
walls the modes are labelled by non-negative integers $(m,n,p)$,
\begin{equation}
\vk=\Big(\frac{m\pi}a,\frac{n\pi}b,\frac{p\pi}d\Big),\quad
\begin{aligned}
E_x&=A_x\cos k_xx\,\sin k_yy\,\sin k_zz,\\
E_y&=A_y\sin k_xx\,\cos k_yy\,\sin k_zz,\\
E_z&=A_z\sin k_xx\,\sin k_yy\,\cos k_zz,
\end{aligned}
\quad \omega=c|\vk| .
\end{equation}
The tangential components vanish on the walls automatically, and
$\divg\vE=0$ is $\vk\cdot\mathbf A=0$.  If $m,n,p\ge1$ there are two
independent polarizations; if exactly one index is zero there is one
(the components that carry the vanishing index die); if two are zero
there is none.  The one-polarization modes on the coordinate planes compensate the
surface term of the separate Dirichlet and Neumann counts, so that for a
rectangular box the first correction to Weyl's law is linear in $\omega$
(proportional to the total edge length) rather than quadratic.
\end{example}

\subsection{Waveguides}

For a cylinder (cross-section $\Sigma$) with perfectly conducting walls,
fields $\propto e^{\ii(k_zz-\omega t)}$ split into TM waves
($B_z=0$, $E_z=\psi$ with $-\nabla_\perp^2\psi=\gamma^2\psi$, $\psi|_{\pa\Sigma}=0$) and TE waves
($E_z=0$, $B_z=\psi$ with $\pa_n\psi|_{\pa\Sigma}=0$), with
\begin{equation}
\omega^2=c^2(k_z^2+\gamma^2)
\end{equation}
and a cut-off frequency $c\gamma$ below which the wave does not propagate.  When
$\Sigma$ is \emph{multiply connected} (coaxial line) there is in
addition a TEM wave, $\omega=ck_z$, $\vE_\perp=-\grad_\perp\phi$ with
$\phi$ constant on each boundary component: the waveguide version of the
topological static sector above.

\subsection{Exterior problems and the radiation condition}

For time-harmonic fields $e^{-\ii\omega t}$ and $k=\omega/c$ the equations
become $\curl\vE=\ii k\vB$, $\curl\vB=-\ii k\vE+4\pi\vj/c$, so that
$\curl\curl\vE-k^2\vE=(4\pi\ii k/c)\vj$.  For an exterior domain (a
scatterer, an antenna) the problem
is well posed only with a condition at infinity that selects the
\emph{outgoing} solution, the time-harmonic version of the
retarded Green function.  This is the Silver--M\"uller radiation condition
\begin{equation}\label{eq:SM}
\lim_{r\to\infty}r\,\big(\vB\times\hat{\mathbf r}-\vE\big)=0
\end{equation}
(for an outgoing wave $\vB=\hat{\mathbf r}\times\vE$ and the bracket vanishes at leading order).
For a perfectly conducting scatterer $D$ with incident field
$(\vE^{\rm inc},\vB^{\rm inc})$, one seeks the scattered field with
\begin{equation}
\curl\curl\vE^{\rm sc}-k^2\vE^{\rm sc}=0\ \text{in }\mathbb R^3\setminus\overline D,\quad
\vn\times(\vE^{\rm sc}+\vE^{\rm inc})=0\ \text{on }\pa D,\quad\text{and \eqref{eq:SM}.}
\end{equation}
For real $k>0$ this problem has exactly one solution (Rellich's lemma
and unique continuation).  Far away, $\vE^{\rm sc}\simeq r^{-1}e^{\ii kr}\vE_\infty(\hat{\mathbf r})$
with $\vE_\infty\perp\hat{\mathbf r}$, the \emph{far-field pattern},
from which the scattering cross-section is computed.  The condition
\eqref{eq:SM} is a statement about \emph{causality}: the retarded Green
function is the boundary value as $\omega\to\omega+\ii0$ of a function
analytic in $\operatorname{Im}\omega>0$, and the
analyticity in the upper half plane is equivalent to the vanishing of
the response before the cause (Kramers--Kronig-type relations).

\subsection{Static boundary-value problems}

For time-independent fields, \eqref{eq:maxwell-3d} decouples into
electrostatics, $\curl\vE=0$, $\divg\vE=4\pi\rho$, i.e.\ $\vE=-\grad\varphi$,
$\nabla^2\varphi=-4\pi\rho$, and magnetostatics, $\divg\vB=0$,
$\curl\vB=4\pi\vj/c$ (which requires $\divg\vj=0$).  The standard
uniqueness theorems follow from the energy identity
$\int|\grad(\varphi_1-\varphi_2)|^2\,\dd^3x=\oint(\varphi_1-\varphi_2)\,\pa_n(\varphi_1-\varphi_2)\,\dd a$:
the Dirichlet problem has a unique solution, the Neumann problem is unique
up to a constant and solvable only if $\oint\pa_n\varphi\,\dd a=-4\pi Q$,
and in the presence of conductors the field is unique when either the
potential or the total charge of each conductor is specified.  In a
non-simply-connected region the static magnetic problem requires,
in addition, the circulation of $\vB$ around each independent loop,
again the topological data of the previous subsection.

\subsection{Other problems}

\paragraph{Characteristic initial-value problem.}  Data may also be given on a null
surface instead of a spacelike one.  On a null hyperplane $x^0+x^3=0$ the
free data are only the \emph{two} transverse components of the potential in
light-cone gauge, as functions of the three variables on the hyperplane:
two functions instead of the four on a spacelike slice.  There is no
contradiction: on a null surface the $(q,p)$ pair collapses into a
single function, since the two are related by a constraint
that is first order along the null generators.  Both
describe the same two degrees of freedom.

\paragraph{Data on a timelike surface: ill-posedness.}  Given $\vE,\vB$ and
their normal derivatives on the plane $z=0$ for all $t$ (the ``sideways'' problem, as in
reconstructing a source from measured near fields), the problem is
\emph{ill posed} in the sense of Hadamard.  For a transverse wave
vector $\vk_\perp$ with $|\vk_\perp|>\omega/c$ the waves are
evanescent, decaying as $\exp(-\sqrt{k_\perp^2-\omega^2/c^2}\,|z|)$
away from the source; in the reverse direction they grow
exponentially, so errors in the high-$|\vk_\perp|$ part of the data are
amplified without bound.  This is why the resolution of back-propagated
images is limited to a fraction of the wavelength.

\paragraph{The inverse source problem: non-uniqueness.}  Even with the field
known exactly outside a bounded region the source is not unique.  For
time-harmonic fields choose any smooth compactly supported vector field
$\mathbf F$ and set
$\vj=\dfrac{c^2}{4\pi\ii\omega}\big(\curl\curl\mathbf F-k^2\mathbf F\big)$:
the solution of \eqref{eq:maxwell-3d} is $\vE=\mathbf F$, which vanishes outside the
support, so the source is \emph{non-radiating}.  In the time-dependent
case any spherically symmetric charge distribution with the same
total charge produces the same exterior field (the Coulomb field,
$\vB=0$).  A measurement of the exterior field determines the
source only modulo non-radiating sources.

% ==========================================================================
\section{Discretization and the preservation of constraints}\label{sec:mx-num}
% ==========================================================================

The Cauchy problem is the prototype of the numerical treatment of
Maxwell's equations, and it shows in a particularly clean way what the
constraints are.  On a rectangular lattice (Yee) $\vE$ is stored on edges
and $\vB$ on faces; $\vE$ is known at integer time steps $t^n=n\,\Delta t$ and
$\vB$ at half-integer ones (leapfrog).  Let $\mathcal C$ be the discrete
curl from edge values to face values, $\mathcal C^{\mathsf T}$ its transpose
(face to edge, the curl from faces back to edges), and $\mathcal D$, $\mathcal D'$
the discrete divergences of $\vB$ (on faces) and of $\vE$ (on edges).
The scheme is
\begin{equation}\label{eq:yee}
\begin{aligned}
\vB^{n+1/2}&=\vB^{n-1/2}-c\,\Delta t\ \mathcal C\,\vE^n,\\
\vE^{n+1}&=\vE^n+c\,\Delta t\ \mathcal C^{\mathsf T}\vB^{n+1/2}-4\pi\,\Delta t\ \vj^{n+1/2}.
\end{aligned}
\end{equation}
The staggering is arranged so that the \emph{lattice identities}
\begin{equation}
\mathcal D\,\mathcal C=0,\qquad \mathcal D'\,\mathcal C^{\mathsf T}=0
\end{equation}
hold exactly; these are the lattice versions of $\divg\curl\equiv0$ and
of $\dd^2=0$.  Consequently:
\begin{enumerate}
\item $\mathcal D\vB^{n+1/2}=\mathcal D\vB^{n-1/2}$: the magnetic constraint
  $\divg\vB=0$ is preserved \emph{exactly} (up to round-off), step by step,
  exactly as \eqref{eq:propagate-B} states for the continuum.
\item $\mathcal D'\vE^{n+1}-4\pi\rho^{n+1}=\mathcal D'\vE^n-4\pi\rho^n$ if
  and only if the discrete continuity equation
  $(\rho^{n+1}-\rho^n)/\Delta t+\mathcal D'\vj^{n+1/2}=0$ holds.  The current must be
  deposited on the lattice in a charge-conserving way, otherwise
  Gauss's law is violated from the first step on, which is the discrete
  form of the compatibility condition \eqref{eq:continuity}.
\item For $\vj=0$ the quadratic form
  $W^n=\tfrac12|\vE^n|^2+\tfrac12\,\vB^{n-1/2}\!\cdot\vB^{n+1/2}$ (with the usual lattice
  inner products, and up to the factor $1/4\pi$) is conserved exactly:
  $W^{n+1}=W^n$: the contribution of the $\vE$ update,
  $+\tfrac12c\,\Delta t\,\vB^{n+1/2}\!\cdot\mathcal C(\vE^{n+1}+\vE^n)$, cancels that of the
  $\vB$ update, which is the same expression with the opposite sign.
  $W^n$ is positive definite under the stability condition below.
\item Substituting plane waves gives the numerical dispersion relation
  \begin{equation}
  \frac{\sin^2(\omega\Delta t/2)}{(c\Delta t)^2}=\sum_{i=1}^3\frac{\sin^2(k_i\Delta x_i/2)}{\Delta x_i^2},
  \end{equation}
  so the scheme is second-order accurate, the phase speed is anisotropic and
  smaller than $c$ at short wavelengths, and the scheme is stable if and only if
  the Courant--Friedrichs--Lewy condition
  $c\,\Delta t\le\big(\sum_i\Delta x_i^{-2}\big)^{-1/2}$ holds ($c\Delta t\le\Delta x/\sqrt3$
  for a cubic lattice), the discrete counterpart of the statement that the
  numerical domain of dependence must contain the physical one.
\end{enumerate}
The lesson is general: a discretization that respects the structure
\emph{constraints preserved by the evolution} reproduces the geometry of the
continuum theory exactly, whereas a discretization that merely approximates
each term does not, and the violation of the constraints grows.  Note also what kind of numerics this is: the lattice computation
approximates a quantity, the field, that already has a real-world
referent, in contrast to lattice formulations whose computed object has
none.

% ==========================================================================
\section{Summary}\label{sec:mx-summary}
% ==========================================================================

\begin{itemize}
\item The vacuum Maxwell equations are $8$ equations for $6$ unknowns,
  with one identity and one compatibility condition (charge
  conservation): $6$ evolution equations and $2$ constraints which the
  evolution preserves.
\item The field has $2$ degrees of freedom per point, i.e.\ $4$ free
  functions of three variables: the transverse parts of $\vE$ and $\vB$.
  The longitudinal parts are slaved to the instantaneous charge density
  (Gauss's law) or vanish; gauge freedom removes the rest of the
  potential.
\item Geometrically, $F=\dd A$ is a closed $2$-form, $\dd{\star F}=(4\pi/c)\star J$;
  the theory is Poincar\'e-, conformally- and (without sources)
  duality-invariant, $F_{\mu\nu}$ has two invariants and a null/non-null classification,
  and $\mathbf F=\vE+\ii\vB$ transforms under $SO(3,\mathbb C)$.
\item The Cauchy problem is well posed (symmetric hyperbolic); the
  maximal speed is $c$; free-space solutions obey the strong Huygens
  principle.
\item Retarded and advanced solutions are both solutions; the
  distinction between them is a condition at infinity, not part of the
  local field equations.
\item In bounded domains the perfect-conductor problem is conservative, has a
  discrete spectrum of cavity modes, and a topological static sector;
  exterior problems require a radiation condition; problems with data on
  timelike surfaces are ill posed; inverse source problems are not unique.
\item A discretization that mirrors $\dd^2=0$ preserves the
  constraints exactly.
\end{itemize}

\subsection*{Problems}

\begin{problem}
Derive \eqref{eq:wave-EB} from \eqref{eq:maxwell-3d}.  Show that
$\Box\vE=0=\Box\vB$ does \emph{not} imply \eqref{eq:maxwell-3d} (count
the Cauchy data).
\end{problem}

\begin{problem}
Using \eqref{eq:boost}, show that if $\vE\cdot\vB=0$ and $|\vE|<|\vB|$ then
in the frame moving with velocity $\vv=c\,\vE\times\vB/\vB^2$ the
electric field vanishes.  What happens when $|\vE|=|\vB|$?
\end{problem}

\begin{problem}
Verify that $F_{\mu\nu}F^{\mu\nu}=2(\vB^2-\vE^2)$,
$F_{\mu\nu}\Ftil^{\mu\nu}=-4\vE\cdot\vB$, and that
$\det(\lambda-F^\mu{}_\nu)=\lambda^4-I_1\lambda^2-I_2^2$ with
$I_1=\vE^2-\vB^2$, $I_2=\vE\cdot\vB$.
\end{problem}

\begin{problem}
Verify $\pa_\mu T^{\mu\nu}=-c^{-1}F^{\nu\lambda}j_\lambda$ from
\eqref{eq:T} and \eqref{eq:maxwell-cov}, and show that $T^\mu{}_\mu=0$.
\end{problem}

\begin{problem}
Show that in $D$ spacetime dimensions the Maxwell field has $D-2$ degrees of
freedom, by the Hamiltonian counting of Section~\ref{sec:mx-dof}
and by the little-group argument.
\end{problem}

\begin{problem}
Prove Proposition~\ref{prop:local-energy} directly from the Kirchhoff
formula \eqref{eq:kirchhoff}.  Show that the strong Huygens principle fails
for waves independent of $z$ (hint: integrate the Green function over $z$).
\end{problem}

\begin{problem}
In the rectangular cavity, count the modes below a given $\omega$ for a cube of side $a$
and compare with Weyl's law $V\omega^3/(3\pi^2c^3)$; find the lowest resonance
and its degeneracy.
\end{problem}

\begin{problem}
Show that the quantity $W^n$ of Section~\ref{sec:mx-num} is conserved by the scheme
\eqref{eq:yee} for $\vj=0$, using only
that $\mathcal C^{\mathsf T}$ is the transpose of $\mathcal C$.  Show that a
current deposition violating discrete continuity produces a Gauss-law
residual that grows linearly in time for a steady error.
\end{problem}

\begin{problem}
A spherical shell carries a time-dependent radial polarization
$\mathbf P(r,t)\hat{\mathbf r}$, with $\rho=-\divg\mathbf P$ and
$\vj=\pa_t\mathbf P$.  Show that the exterior field vanishes, and explain
why the spherically symmetric part of any source never radiates.
\end{problem}
